\documentclass[14pt,compress,aspectratio=169]{beamer} % selection of font size = 9, 10, 11, 12, 14, ....pt


% specify the language of the document here
%
\newcommand\setslidelanguage{english}
%\newcommand\setslidelanguage{german}

%
% this flag turns on a basic Fraunhofer CI overwrite
% this just changes colors atm, types and more convenient
% theme handling have to be added
%
%\newcommand\setslidetheme{tui} % - alpha
%\newcommand\setslidetheme{tuiiis} % - alpha
\newcommand\setslidetheme{tuiizfp} % - alpha
%\newcommand\setslidetheme{iis} % - alpha
%\newcommand\setslidetheme{iistui} % - alpha
%\newcommand\setslidetheme{izfp} % - alpha
%\newcommand\setslidetheme{izfptui} % - alpha


%
% load header with all includes
%
%%% input encoding, font encoding
\usepackage[T1]{fontenc}
\usepackage[utf8]{inputenc}
\usepackage{lmodern}


%%% beamerstuff - requirements
\usepackage{ifthen}
\usepackage[absolute,overlay]{textpos}
\setlength{\TPHorizModule}{1cm}
\setlength{\TPVertModule}{1cm}


%%% presentation style stuff
% load a lot of custom macros and sets the
% beamer theme
% when new logos and stuff are required
% look there
%%% presentation style stuff


%%%%%%%%%%%%%%%%%%%%%%%%%%%%%%%%%%%%%%%%%%%%%%%%%%%%%%%%%%%%%%%%%%%%%%%%%%%%%%%%%%%%%%%%%%%%%%%%%%%%%%%%%%%%%%%%%%%%%%%%%%%%%%%%%%%%%
% general stuff

% thanks to https://tex.stackexchange.com/questions/2541/beamer-frame-numbering-in-appendix
% this helps us to organize framenumbering s.t. they do not increase for the appendix
\newcommand{\backupbegin}{
   \newcounter{framenumberappendix}
   \setcounter{framenumberappendix}{\value{framenumber}}
}
\newcommand{\backupend}{
   \addtocounter{framenumberappendix}{-\value{framenumber}}
   \addtocounter{framenumber}{\value{framenumberappendix}} 
}


% this is required to omit problems with tikz-externalize
% as I did all that footer/header stuff using tikz
% so we have to prevent tikz externalize from
% externalizing these
% otherwise framecounters and sections will screw up
\newcommand{\donotexternalize}[1]{
\tikzset{external/export=false}
#1
\tikzset{external/export=true}
}


%%%%%%%%%%%%%%%%%%%%%%%%%%%%%%%%%%%%%%%%%%%%%%%%%%%%%%%%%%%%%%%%%%%%%%%%%%%%%%%%%%%%%%%%%%%%%%%%%%%%%%%%%%%%%%%%%%%%%%%%%%%%%%%%%%%%%
% TU Ilmenau stuff

%% Logos
\newcommand\logotuiwhite{\includegraphics[width=1.86cm]{logos/tu_bw_white.pdf}}
\newcommand\logotuicolor{\includegraphics[width=1.86cm]{logos/tu_col_green.pdf}}

\newcommand\settuitextmacros{
	% Sets macros for text formatting
	\newcommand<>{\cihighlight}[1]{\textcolor##2{tui_orange_dark}{##1}}
	\newcommand<>{\cigray}[1]{\textcolor##2{tui_gray_dark}{##1}}
	\newcommand<>{\cihide}[1]{\textcolor##2{white}{##1}}
}


%% Style
\newcommand{\setTUIStyleTwoLogos}{
	\mode<presentation>
		{
			\usepackage{theme/beamerthemetui}
			\setbeamercovered{transparent}
		}
	}
	
	
	
% for enabling 16:9
\newlength{\ruleswidth}
\setlength{\ruleswidth}{\dimexpr(\paperwidth-20mm)\relax}
	
	
%%%%%%%%%%%%%%%%%%%%%%%%%%%%%%%%%%%%%%%%%%%%%%%%%%%%%%%%%%%%%%%%%%%%%%%%%%%%%%%%%%%%%%%%%%%%%%%%%%%%%%%%%%%%%%%%%%%%%%%%%%%%%%%%%%%%%
% Fraunhofer stuff


%% Logos
\newcommand\logoizfpwhite{\raisebox{-6.68mm}{\includegraphics[width=1.86cm]{logos/izfp_85mm_p334_w.pdf}}}
\newcommand\logoizfpcolor{\raisebox{-6.68mm}{\includegraphics[width=1.86cm]{logos/izfp_85mm_p334-eps-converted-to.pdf}}}
\newcommand\logoiiswhite{\raisebox{-6.72mm}{\includegraphics[width=1.86cm]{logos/iis_85mm_p334.pdf}}}
\newcommand\logoiiscolor{\raisebox{-6.72mm}{\includegraphics[width=1.86cm]{logos/iis_85mm_p334.pdf}}}


\newcommand\setfhitextmacros{
	% Sets macros for text formatting
	\newcommand<>{\cihighlight}[1]{\textcolor##2{fraunhofer_orange}{##1}}
	\newcommand<>{\cigray}[1]{\textcolor##2{fraunhofer_silver!70!black}{##1}}
	\newcommand<>{\cihide}[1]{\textcolor##2{white}{##1}}
	% taken from https://tex.stackexchange.com/questions/23034/how-to-get-larger-item-symbols-for-some-lists-in-a-beamer-	
	% presentation and
	% https://tex.stackexchange.com/questions/87133/changing-the-color-of-itemize-item-in-beamer
	\setbeamertemplate{itemize subsubitem}{\tiny\raise0.5pt\hbox{\color{fraunhofer_silver}$ \blacksquare$}}
}

\newcommand{\setFHIStyleTwoLogos}{
	\mode<presentation>
		{
			\usepackage{theme/beamerthemefhi}
			\setbeamercovered{transparent}
		}
	% Big TODO: Frutiger font
}




%%%%%%%%%%%%%%%%%%%%%%%%%%%%%%%%%%%%%%%%%%%%%%%%%%%%%%%%%%%%%%%%%%%%%%%%%%%%%%%%%%%%%%%%%%%%%%%%%%%%%%%%%%%%%%%%%%%%%%%%%%%%%%%%%%%%%
% set everything up using the macros
%% TUI
\ifthenelse{\equal{\setslidetheme}{tui}}
{
		\newcommand\setfirstlogo{\logotuiwhite}
		\newcommand\setsecondlogo{}
		
		\setTUIStyleTwoLogos
		
		\settuitextmacros
}{}
\ifthenelse{\equal{\setslidetheme}{tuiiis}}
{
		\newcommand\setfirstlogo{\logotuiwhite}
		\newcommand\setsecondlogo{\logoiiswhite}
		
		\setTUIStyleTwoLogos
		
		\settuitextmacros
}{}
\ifthenelse{\equal{\setslidetheme}{tuiizfp}}
{
		\newcommand\setfirstlogo{\logotuiwhite}
		\newcommand\setsecondlogo{\logoizfpwhite}
		
		\setTUIStyleTwoLogos
		
		\settuitextmacros
}{}


%% Fraunhofer 
\ifthenelse{ \equal{\setslidetheme}{iis}}
{
		\newcommand\setfirstlogo{\logoiiscolor}
		\newcommand\setsecondlogo{}

		\setFHIStyleTwoLogos
		
		\setfhitextmacros
}{}
\ifthenelse{ \equal{\setslidetheme}{izfp}}
{
		\newcommand\setfirstlogo{\logoizfpcolor}
		\newcommand\setsecondlogo{}
		
		\setFHIStyleTwoLogos
		
		\setfhitextmacros
}{}
\ifthenelse{ \equal{\setslidetheme}{iistui}}
{
		\newcommand\setfirstlogo{\logoiiscolor}
		\newcommand\setsecondlogo{\logotuicolor}

		\setFHIStyleTwoLogos
		
		\setfhitextmacros
}{}
\ifthenelse{ \equal{\setslidetheme}{izfptui}}
{
		\newcommand\setfirstlogo{\logoizfpcolor}
		\newcommand\setsecondlogo{\logotuicolor}

		\setFHIStyleTwoLogos
		
		\setfhitextmacros
}{}




%%% language settings
\ifthenelse{ \equal{\setslidelanguage}{german}}
	{	\usepackage[ngerman]{babel} }
	{	\usepackage[english]{babel} }


%%% general (useful) packages
% math
\usepackage{amsmath}
\usepackage{bm}
\usepackage{mathtools} % For \coloneqq

% graphics
\usepackage{pgfplots}
\pgfplotsset{compat=newest}
\usepackage{animate}
\usepackage{tikz}
\usepackage{multido}
\usetikzlibrary{patterns}
\usetikzlibrary{matrix,arrows,shapes,trees,fit, decorations, spy}
\usetikzlibrary{calc,positioning,scopes,backgrounds}
\usetikzlibrary{tikzmark,decorations.pathreplacing,calligraphy}
\usepgfplotslibrary{colormaps}
\usetikzlibrary{external} % most useful, but requires --shell-escape
\usetikzlibrary{petri} %for tokens in the block diagram for iterative GD
% this can be used if you want to use tikzexternalize
% it precompiles your plots to a subfolder called
% figures/tikz-ext/ and names them according to
% the section 
% see https://tex.stackexchange.com/questions/177292/chapter-and-figure-numbers-in-tikz-externalize-prefix
%
%
%\tikzexternalize[
%    prefix=figures/tikz-ext/,
%    figure name=plot_sec\thesubsection_no,
%]

%% for putting figures at the specific locations
\usepackage[absolute,overlay]{textpos}
\setlength{\TPHorizModule}{1mm}
\setlength{\TPVertModule}{1mm}

% for checkmark in conclusion
\usepackage{bbding}


% settings and default macros
\usepackage[load-configurations=abbreviations,load-configurations=binary,binary-units=true]{siunitx}	%korrektes Setzen von SI-Einheiten
\DeclareSIUnit{\mms}{\milli\squaremetre}
\DeclareSIUnit{\inch}{in}
\DeclareSIUnit{\inchs}{in\squared}
\DeclareSIUnit{\mil}{mil}
\DeclareSIUnit{\Msps}{Msps}
\DeclareSIUnit{\Mbps}{Mbps}
\DeclareSIUnit{\LSB}{LSB}
\DeclareSIUnit{\pFS}{\percent FS}
\DeclareSIUnit{\dBc}{\deci\bel c}
\DeclareSIUnit{\dBm}{\deci\bel m}
\DeclareSIUnit{\dBFS}{\deci\bel FS}
\DeclareSIUnit{\dB}{\deci\bel}
\DeclareSIUnit{\dBi}{\deci\bel i}
\DeclareSIUnit{\hex}{0x}
\DeclareSIUnit{\vp}{\volt_{\text{p}}}
\DeclareSIUnit{\vpp}{\volt_{\text{pp}}}
\DeclareSIUnit{\kb}{\kilo\bit}
\DeclareSIUnit{\kB}{\kilo\byte}
\DeclareSIUnit{\MB}{\mega\byte}




%
% your settings and includes go here
%
%%%% bibliography %%%%%
\usepackage{url}
\usepackage{natbib} % required for the citation example
\usepackage{bibentry} % required for the citation example
%
\nobibliography* % required for the citation example
\bibliographystyle{plain}

% "newcommand" collections
% important!!!! no numbers in the name of commands!!!!!!!!

%=====================  for math =======================%
% symbols
% scan positions: p
\newcommand{\pp}{\bm{p}}
\newcommand{\pphat}{\bm{\hat{p}}}
\newcommand{\ppdelta}{\Delta \pp}
\newcommand{\ppdeltahat}{\Delta \hat{\pp}}
% scan positions: x
\newcommand{\xdelta}{\Delta x}
\newcommand{\xhat}{\hat{x}}
\newcommand{\xopt}{\hat{x}_{\optimized}}
\newcommand{\xdeltaest}{\xdelta_{\estimated}}
\newcommand{\xdeltaopt}{\xdelta_{\optimized}}
\newcommand{\xdeltahat}{\hat{\xdelta}}
\newcommand{\xvec}{\bm{x}}
\newcommand{\xvechat}{\bm{\hat{x}}}
\newcommand{\xdeltavec}{\Delta \xvec}
\newcommand{\xtrackvec}{\xvechat_{\track}}
\newcommand{\xoptvec}{\xvechat_{\optimized}}
\newcommand{\xdeltaoptvec}{\Delta \xvechat_{\optimized}}
% scatter position
\newcommand{\scatterer}{\bm{s}}
% A-Scans
\newcommand{\ascan}{a}
\newcommand{\ascanvec}{\bm{a}}
\newcommand{\ascank}{\bm{a}_{k}}
\newcommand{\ascanvechat}{\bm{\hat{a}}}
\newcommand{\adothat}{\dot{\bm{\hat{a}}}}
\newcommand{\aopt}{\bm{\tilde{a}}}
\newcommand{\Ascan}{\bm{A}}
% Pulse
\newcommand{\pulse}{h}
\newcommand{\pulsevec}{\bm{h}}
\newcommand{\pulsevecdot}{\bm{h}'}
\newcommand{\SAFTcol}[2]{\pulsevec_{#1}^{( #2)}}
% SAFT matrix
\newcommand{\SAFT}{\bm{H}}
\newcommand{\SAFTp}{\SAFT (\pp)}
\newcommand{\SAFTx}{\SAFT (x)}
\newcommand{\SAFTk}{\SAFT_{k}}
\newcommand{\SAFThat}{\hat{\bm{H}}}
\newcommand{\SAFThatk}{\hat{\bm{H}}_{k}}
\newcommand{\SAFTapproxxvec}{\left( \SAFThat + \bm{E} \Jacobianhat \right)}

% Jacobian relevent
\newcommand{\Jacobian}{\bm{J}} 
\newcommand{\Jacobianhat}{\hat{\Jacobian}} 
\newcommand{\Jacobianpartial}{\bm{J}_{\SAFTcol}}
\newcommand{\Deriv}{\bm{D}} 
\newcommand{\Derivcol}{\bm{d}_{l}} 
% else
\newcommand{\reflectivity}{r}
\newcommand{\refcoeff}{\beta}
\newcommand{\defect}{\bm{b}} % should be modified? 
\newcommand{\defectsingle}{b} % should be modified? 
\newcommand{\defecthat}{\hat{\bm{b}}} % should be modified? 
\newcommand{\zvec}{\bm{z}} 
\newcommand{\noisevec}{\bm{n}} 
\newcommand{\fhatpartial}{\hat{f}_{i}}
\newcommand{\taux}{\tau_{l} (x)}
\newcommand{\Identity}{\bm{I}} 
\newcommand{\SEdag}{\SE^{\dagger}} 
\newcommand{\MSEdag}{\MSE^{\dagger}} 

% Math functions/operators
\newcommand{\norm}[1]{%
	\left\lVert#1\right\rVert
}
\newcommand{\vecoperator}[1]{%
	\vectorize \{ #1 \}
}
\newcommand{\diagoperator}[1]{%  
     \operatorname{diag} \{ #1 \}
}
\newcommand{\Real}{%  
     \operatorname{Re}
}
\newcommand{\Imag}{%  
     \operatorname{Im}
}



% some special characters
\newcommand{\ii}{\mathrm{i}}
\newcommand{\dd}{\mathrm{d}}
\newcommand{\ee}{\mathrm{e}}
\newcommand{\dirac}{\delta}

\newcommand{\RR}{{\mathbb{R}}}
\newcommand{\NN}{{\mathbb{N}}}
\newcommand{\CC}{{\mathbb{C}}}
\newcommand{\OO}{{\mathcal{O}}}


%%%%%%%%%%%%%%%%%%%%%%%%%%%%%%%%%%%%%%%%%%%%%%
%======================================== table macro ====%
%%%%%%%%%%%%%%%%%%%%%%%%%%%%%%%%%%%%%%%%%%%%%%
% table scaling
\newcommand{\inputTable}[2]{%  
	\resizebox{#1}{!}{%
	 \input{#2}
	}
}


%%%%%%%%%%%%%%%%%%%%%%%%%%%%%%%%%%%%%%%%%%%%%%
%========================================== for TikZ  ====%
%%%%%%%%%%%%%%%%%%%%%%%%%%%%%%%%%%%%%%%%%%%%%%
\newcommand{\Nx}{3}
\newcommand{\Ny}{2}
\newcommand{\ddx}{1.5cm}
\newcommand{\ddy}{\ddx}
\newcommand{\ddz}{1.5cm}
\newcommand{\rCircle}{0.11cm}
\newcommand{\rCircleCamera}{0.07cm}

%%% draw transducer %%%
%5 3D
\newcommand{\drawTransducer}[4]{ % scale,scaley, x, y in axis
	\draw[draw = none, fill = white] (axis cs: #3 - #1, #4 - #2, 0) -- (axis cs: #3 + #1, #4 - #2, 0) -- (axis cs: #3 + #1, #4 + #2, 0) --  (axis cs: #3 + #1, #4 + #2, -#1) -- (axis cs: #3 - #1, #4 + #2, -#1) -- (axis cs: #3 - #1, #4 - #2, -#1) -- (axis cs: #3 - #1, #4 - #2, 0);
	\draw[] (axis cs: #3 - #1, #4 - #2, 0) -- (axis cs: #3 + #1, #4 - #2, 0) -- (axis cs: #3 + #1, #4 - #2, -#1) -- (axis cs: #3 - #1, #4 - #2, -#1) -- (axis cs: #3 - #1, #4 - #2, 0);
	\draw[] (axis cs: #3 + #1, #4 - #2, 0) -- (axis cs: #3 + #1, #4 + #2, 0) -- (axis cs: #3 + #1, #4 +  #2, -#1) -- (axis cs: #3 - #1, #4 +  #2, -#1) -- (axis cs: #3 - #1, #4 - #2, -#1);
	\draw[] (axis cs: #3 + #1, #4 - #2, -#1) -- (axis cs: #3 + #1, #4 +  #2, -#1);
}

%% 2D
\newcommand{\drawTransducerTwoD}[2]{ % scale, x  1,3
	\draw[draw = none, fill = white] (axis cs: #2 - #1, 0) -- (axis cs: #2 + #1, 0) --  (axis cs: #2 + #1, -#1) -- (axis cs: #2 - #1, -#1) -- (axis cs: #2 - #1, 0);
	\draw[] (axis cs: #2 - #1, 0) -- (axis cs: #2 + #1, 0) -- (axis cs: #2 + #1, -#1) -- (axis cs: #2 - #1, -#1) -- (axis cs: #2 - #1, 0);
	\draw[] (axis cs: #2 + #1, 0) -- (axis cs: #2 + #1, -#1) -- (axis cs: #2 - #1, -#1);
}

%%% draw video camera %%%
%% 3D
\newcommand{\drawCamera}[6]{ %<rotation origin (x, y, z)>, <start x>, <start y>, <start z> in axis, <camera width [dx]>, <camera height [dz]>
	% base
	\draw[rotate around = {30 : (axis cs: #1)}] (axis cs: #2, #3, #4) rectangle (axis cs: #2 + #5, #3, #4 - #6) ;
	% origin of rotation
	\node[campoint] (rotationorg) at (axis cs: #1) {};
	% "trapezoid" part (tip of the camera)
	\draw[rotate around = {30 : (axis cs: #1)}] (axis cs: #2, #3, #4 - 0.25*#6) -- (axis cs: #2 - 0.5* #6, #3, #4) -- (axis cs: #2 - 0.5* #6, #3, #4-#6) -- (axis cs: #2, #3, #4 - 0.75*#6);	
	% cable
	\draw[rotate around = {30 : (axis cs: #1)}] (axis cs: #2 + #5, #3, #4 - 0.5*  #6) .. controls (axis cs: #2 + 1.5* #5, #3, #4 - 0.5*  #6) and (axis cs: #2 + #5, #3, - 0.5*  #6) .. (axis cs: #2 + 1.5* #5, #3, - #6);	
}

%% 2D
\newcommand{\drawCameraTwoD}[5]{ %<rotation origin (x, z)>, <start x>, <start y>, <start z> in axis, <camera width [dx]>, <camera height [dz]> 1, 2, 4, 5, 6 -> 4-> 3, 5->4, 6-> 5
	% base
	\draw[rotate around = {30 : (axis cs: #1)}] (axis cs: #2, #3) rectangle (axis cs: #2 + #4, #3 - #5) ;
	% origin of rotation
	\node[campoint] (rotationorg) at (axis cs: #1) {};
	% "trapezoid" part (tip of the camera)
	\draw[rotate around = {30 : (axis cs: #1)}] (axis cs: #2, #3 - 0.25*#5) -- (axis cs: #2 - 0.5* #5, #3) -- (axis cs: #2 - 0.5* #5, #3-#5) -- (axis cs: #2, #3 - 0.75*#5);	
	% cable
	\draw[rotate around = {30 : (axis cs: #1)}] (axis cs: #2 + #4, #3 - 0.5*  #5) .. controls (axis cs: #2 + 1.5* #4, #3 - 0.5*  #5) and (axis cs: #2 + #4, - 0.5*  #5) .. (axis cs: #2 + 1.5* #4, - #5);	
}


%%%%%%%%%%%%%%%%%%%%%%%%%%%%%%%%%%%%%%%%%%%%%%
%======================================= image macro ====%
%%%%%%%%%%%%%%%%%%%%%%%%%%%%%%%%%%%%%%%%%%%%%%
% TikZ scaling
\newcommand{\inputTikZ}[2]{%<scaling factor>, <name of the tex file>
     \scalebox{#1}{\input{#2}}  
}

% BRP: for error correction animation
\newcommand{\measanimate}[5]{ % <scale size>, <slide page for base pulse>, <slide page for highlighting the pulse 1mm away>, <slide page for highlighting the pulse 2.5mm away>, <slide page for highlighting the pulse 5mm away>
	\scalebox{#1}{
		\begin{tikzpicture}
			\begin{axis}[
					width=6cm, height=6.3cm,  at={(0.7cm,0.3cm)},
					ticks=none, axis lines = center, 
					xmin=-0.5, xmax=10.5, ymin=-0.45, ymax=10,
					xlabel={$x$}, ylabel={$t$},
					y dir=reverse,
					x label style={at={(axis cs: 10.5, 0)}, anchor= west},
					y label style={at={(axis cs: -0.6, 10)}, anchor = north},
			        ]
			        
			        % Defect
			        \node at (axis cs: 5, 4) {\pgftext{\includegraphics[scale=0.15]{images/defect}}};
					% Pulse
					\only<#2>{
					\input{figures/pytikz/1D/coordinates/pulse/pulse_1_blue.tex}
					\input{figures/pytikz/1D/coordinates/pulse/pulse_2_blue.tex}
					\input{figures/pytikz/1D/coordinates/pulse/pulse_3_blue.tex}
					\input{figures/pytikz/1D/coordinates/pulse/pulse_4_blue.tex}
					\input{figures/pytikz/1D/coordinates/pulse/pulse_5_blue.tex}
					\input{figures/pytikz/1D/coordinates/pulse/pulse_6_blue.tex}
					\input{figures/pytikz/1D/coordinates/pulse/pulse_7_blue.tex}
					\input{figures/pytikz/1D/coordinates/pulse/pulse_8_blue.tex}
					\input{figures/pytikz/1D/coordinates/pulse/pulse_9_blue.tex}
					}
					
					% Highlight
					\only<#3>{\input{figures/pytikz/1D/coordinates/pulse/pulse_4_highlight.tex}} %1 mm away
					\only<#4>{\input{figures/pytikz/1D/coordinates/pulse/pulse_3_highlight.tex}} % 2.5mm away
					\only<#5>{\input{figures/pytikz/1D/coordinates/pulse/pulse_1_highlight.tex}} % 5mm away
			
			\end{axis}			
		\end{tikzpicture}
	}
}

\newcommand{\curvefit}[5]{ % <scale size>, <slide page for base pulse>, <slide page for highlighting the pulse 2.5mm away>, <slide page for adding the defect position>, <slide page for adding the scan position>, 
	\scalebox{#1}{
		\begin{tikzpicture}
			\begin{axis}[
					width=6cm, height=6.3cm,  at={(0.7cm,0.3cm)},
					ticks=none, axis lines = center, 
					xmin=-1.5, xmax=10.5, ymin=-1.5, ymax=10,
					xlabel={$x$}, ylabel={$t$},
					y dir=reverse,
					x label style={at={(axis cs: 10.5, 0)}, anchor= west},
					y label style={at={(axis cs: -0.6, 10)}, anchor = north},
			        ]
			        
			        % Defect
			        \node at (axis cs: 5, 4) {\pgftext{\includegraphics[scale=0.15]{images/defect}}};
					% Pulse
					\only<#2>{
					\input{figures/pytikz/1D/coordinates/pulse/pulse_1_blue.tex}
					\input{figures/pytikz/1D/coordinates/pulse/pulse_2_blue.tex}
					\input{figures/pytikz/1D/coordinates/pulse/pulse_3_blue.tex}
					\input{figures/pytikz/1D/coordinates/pulse/pulse_4_blue.tex}
					\input{figures/pytikz/1D/coordinates/pulse/pulse_5_blue.tex}
					\input{figures/pytikz/1D/coordinates/pulse/pulse_6_blue.tex}
					\input{figures/pytikz/1D/coordinates/pulse/pulse_7_blue.tex}
					\input{figures/pytikz/1D/coordinates/pulse/pulse_8_blue.tex}
					\input{figures/pytikz/1D/coordinates/pulse/pulse_9_blue.tex}
					}
					
					% Highlight
					\only<#3>{\input{figures/pytikz/1D/coordinates/pulse/pulse_3_highlight.tex}} %2.5mm away
					% Node: defect positions
					\only<#4>{%
						\node at (axis cs: 5, 0) {$\shortmid$};
						\node at (axis cs: 5, -1) {$x_{\dist}$};
						\node at (axis cs: 0, 4) {$-$};
						\node at (axis cs: -1, 4) {$z_{\dist}$};
					}% 
					% Node: i-th scan
					\only<#5>{%
						% Defect
						\node at (axis cs: 3, 0) {$\shortmid$};
						\node at (axis cs: 3, -1) {$x_{k}$};
						\node at (axis cs: 0, 5) {$-$};
						\node at (axis cs: -1, 5) {$z_{k}$};
					}% 
			
			\end{axis}			
		\end{tikzpicture}
	}
}

\newcommand{\TLSanimate}[3]{ % <scale size>,  <slide page for z_def = 762dz>, <slide page for z_def = 1270dz>, 
	\scalebox{#1}{
		\begin{tikzpicture}
			\begin{axis}[
					width=6cm, height=6.3cm,  at={(0.7cm,0.3cm)},
					ticks=none, axis lines = center, 
					xmin=-1.5, xmax=6, ymin=-1.5, ymax=10,
					xlabel={$x$ [\SI{}{\milli \metre}] }, ylabel={$z$ [\SI{}{\milli \metre}] },
					y dir=reverse,
					x label style={at={(axis cs: 5.7, 0)}, anchor= south},
					y label style={at={(axis cs: -0.6, 10)}, anchor = north},
			        ]
			        
					% z_def = 762dz
					\only<#2>{%
						% Defect
						\node at (axis cs: 3, 4) {\pgftext{\includegraphics[scale=0.15]{images/defect}}};
						% Defect position ticks
						\node at (axis cs: 3, 0) {$\shortmid$};
						\node at (axis cs: 3, -1) {$20$};
						\node at (axis cs: 0, 4) {$-$};
						\node at (axis cs: -1, 4) {$30$};
					}% 
					
					% z_def = 1270dz
					\only<#3>{%
						% Defect
						\node at (axis cs: 3, 6.5) {\pgftext{\includegraphics[scale=0.15]{images/defect}}};
						% Defect position ticks
						\node at (axis cs: 3, 0) {$\shortmid$};
						\node at (axis cs: 3, -1) {$20$};
						\node at (axis cs: 0, 6.5) {$-$};
						\node at (axis cs: -1, 6.5) {$50$};
					}% 

			
			\end{axis}			
		\end{tikzpicture}
	}
}


%%%%%%%%%%%%%%%%%%%%%%%%%%%%%%%%%%%%%%%%%%%%%%%%
%===================== 1D Visualization ======================%
%%%%%%%%%%%%%%%%%%%%%%%%%%%%%%%%%%%%%%%%%%%%%%%%

% ME TLS: 762dz vs 1270dz
\newcommand{\meTLS}[7]{ % <scale size>, <label font size>, <tick font size>,  <slide page to pop up the ME 762dz>, <fname for ME 762dz>, <slide page to pop up the ME 1270dz>, <fname for ME 1270dz>
\scalebox{#1}{
	\pgfplotsset{
			xmin = -0.05, xmax=1.05, 
			ymin = -0.005, ymax=0.15, 
			scaled ticks=false % to avoid formtting with 10^-2 in y tick
	}
	\begin{tikzpicture}
            \begin{axis}[
                width = 10cm, 
            	   height = 6.5cm,
            	   grid=both,
    			   grid style={line width=.1pt, draw=gray!20},
                xlabel = {Tracking error $/ \lambda$ \cigray{(\SI{1.26}{\milli \metre})} },
                ylabel = {$\norm{\pp_{\dist} - \pphat_{\dist}}_{2}$ $/ \lambda$},
                label style = {font = #2},
                tick label style = {
                		/pgf/number format/fixed, % to avoid formtting with 10^-2 in y tick
                		font = #3
                },
                %y dir = reverse,
                %xtick = {0, 10, ..., 30}, %to customize the axis
                %xticklabel = {0, 5, 10, 15}
                extra y ticks={0.125},
                %extra y tick labels = {$\SE^{\dagger}_{\threshold}$}, 
                %extra y tick style={font = #2},
                ]
                \only<#4>{
	                	\input{#5}
	                	% y = 1.25lamda line
		             \addplot[gray, dashed, mark = , line width = 2pt] coordinates{
		            			(0.0, 0.1243)
		            			(1.0, 0.1243)
		            };
                }
                \only<#6>{
                		\input{#7}
                		% y = 0.096lamda line
%                		\addplot[gray, dashed, mark = , line width = 2pt] coordinates{
%		            			(0.0, 0.096)
%		            			(1.0, 0.096)
%		            };
                	} 
            \end{axis}
	\end{tikzpicture}
	}
}

%%%%% Result: evaluation
% API  results with animation
\newcommand{\resultAPIanimate}[6]{ %<scale size>, <font size>, <slide page for mark>, <mark coordinate for Reco_true>, <mark coordinate for Reco_track>, <mark coordinate for Reco_opt> 
\scalebox{#1}{
	\begin{tikzpicture}
            \begin{axis}[
                width = 7cm, height = 4cm,
                xlabel = {ROI depth [\SI{}{\milli \metre}]}, ylabel = {$\MAPI$},
                ymin= 16.5, ymax= 28,
                label style = {font = #2},
                tick label style = {font = #2},
                xtick = {20, 30, ..., 80},
                ytick = {18, 20, ..., 28},
                grid=both, grid style={line width=.1pt, draw=gray!20},
                legend style ={
                	at={(1.5, 0.8)},
                	nodes={scale=0.95, transform shape},
                	font = #2
                }
                ]
                \input{figures/pytikz/1D/api_true_depth.tex} % true
                \input{figures/pytikz/1D/api_track_depth.tex} % track
                \input{figures/pytikz/1D/api_opt_depth.tex} % opt
                
             % x = 20mm line
             \only<#3>{
             		% true
		         \addplot[tui_red, mark = star, mark size = 2pt] coordinates{
		          		#4
		            };
		         % track
		         \addplot[tui_red, mark = star , mark size = 2pt] coordinates{
		          		#5
		            };
		       % opt
		         \addplot[tui_red, mark = star , mark size = 2pt] coordinates{
		          		#6
		            };
		    }   
                
             % legend
             % To insert the legend title
		   	%\addlegendimage{empty legend} 
		   	% Legend entries  
		   	\addlegendentry{Reference}
             \addlegendentry{No correction}
             \addlegendentry{BEC}
             % Title
             %\addlegendentry{$|s_{x} - x|$}        
            \end{axis}
	\end{tikzpicture}
	}
}

% SE results 
\newcommand{\resultSE}[7]{ % <scale size>, <font size>, <xlabel>, <ymax>, <xtick>, <fname for track>, <fname for opt>
\scalebox{#1}{
	\begin{tikzpicture}
            \begin{axis}[
                width = 7cm, height = 4cm,
                xlabel = {#3}, ylabel = {$\MSEdag$},
                ymin= -0.04, ymax= #4,
                label style = {font = #2},
                tick label style = {font = #2},
                xtick = {#5},
                ytick = {0, 0.2, ..., #4}, 
                grid=both, grid style={line width=.1pt, draw=gray!20},
                legend style ={
                	at={(1.5, 0.8)},
                	nodes={scale=0.95, transform shape},
                	font = #2
                }
                ]
                \input{#6} % track
                \input{#7} % opt
                
             % legend
             % To insert the legend title
		   	%\addlegendimage{empty legend} 
		   	% Legend entries  
             \addlegendentry{No correction}
             \addlegendentry{BEC} 
             % Title
             %\addlegendentry{$|s_{x} - x|$}        
            \end{axis}
	\end{tikzpicture}
	}
}

% API results
\newcommand{\resultAPI}[9]{ % <scale size>, <font size>, <xlabel>, <ymax>, <xtick>, <ytick>, <fname for true>, <fname for track>, <fname for opt>
\scalebox{#1}{
	\begin{tikzpicture}
            \begin{axis}[
                width = 7cm, height = 4cm,
                xlabel = {#3}, ylabel = {$\MAPI$},
                ymin= 16.5, ymax= #4,
                label style = {font = #2},
                tick label style = {font = #2},
                xtick = {#5},
                ytick = {#6},
                grid=both, grid style={line width=.1pt, draw=gray!20},
                legend style ={
                	at={(1.5, 0.8)},
                	nodes={scale=0.95, transform shape},
                	font = #2
                }
                ]
                \input{#7} % true
                \input{#8} % track
                \input{#9} % opt
                
             % legend
             % To insert the legend title
		   	%\addlegendimage{empty legend} 
		   	% Legend entries  
		   	\addlegendentry{Reference}
             \addlegendentry{No correction}
             \addlegendentry{BEC}
             % Title
             %\addlegendentry{$|s_{x} - x|$}        
            \end{axis}
	\end{tikzpicture}
	}
}


% GCNR results
\newcommand{\resultGCNR}[9]{ % <scale size>, <font size>, <xlabel>, <ymin>, <xtick>, <ytick>, <fname for true>, <fname for track>, <fname for opt>
\scalebox{#1}{
	\begin{tikzpicture}
            \begin{axis}[
                width = 7cm, height = 4cm,
                xlabel = {#3}, ylabel = {$\MGCNR$},
                ymin= #4, ymax= 0.97,
                label style = {font = #2},
                tick label style = {font = #2},
                xtick = {#5},
                ytick = {#6},  % 1.01 = otherwise the tick does not show up 
                grid=both, grid style={line width=.1pt, draw=gray!20},
                legend style ={
                	at={(1.5, 0.8)},
                	nodes={scale=0.95, transform shape},
                	font = #2
                }
                ]
                \input{#7} % true
                \input{#8} % track
                \input{#9} % opt
                
             % legend
             % To insert the legend title
		   	%\addlegendimage{empty legend} 
		   	% Legend entries  
		   	\addlegendentry{Reference}
             \addlegendentry{No correction}
             \addlegendentry{BEC}
             % Title
             %\addlegendentry{$|s_{x} - x|$}        
            \end{axis}
	\end{tikzpicture}
	}
}
%%%%%%%%%%%%%%%%%%%%%%%%%%%%%%%%%%%%%%%%%%%%%%%%
%===================== 2D Visualization ======================%
%%%%%%%%%%%%%%%%%%%%%%%%%%%%%%%%%%%%%%%%%%%%%%%%

% cimg with both x- & y-labels
\newcommand{\cimgbothlabels}[4]{% <scale size>, <label font size>, <tick font size>, <png file name>
\scalebox{#1}{
	\begin{tikzpicture}
            \begin{axis}[
                enlargelimits = false,
                axis on top = true,
                axis equal image,
                point meta min = -1,   
                point meta max = 1,
                xlabel = {$x$ in \SI{}{\milli\meter}},
                ylabel = {$y$ in \SI{}{\milli\meter}},
                label style = {font = #2},
                tick label style = {font = #3},
                %y dir = reverse,
                %xtick = {0, 10, ..., 30}, %to customize the axis
                %xticklabel = {0, 5, 10, 15}
                ]
                \addplot graphics [
                    xmin = 0,
                    xmax = 20,
                    ymin = 0,
                    ymax = 20
                ]{#4};
            \end{axis}
	\end{tikzpicture}
	}
}




%%%%%%%%%%%%%%%%%%%%%%%%%%%%%%%%%%%%%%%%%%%%%%
%=================== 2D visualization ======================%
%%%%%%%%%%%%%%%%%%%%%%%%%%%%%%%%%%%%%%%%%%%%%%
% cimg for z_def = 20mm
\newcommand{\imgzdefshallow}[4]{% <scale size>, <label font size>, <tick font size>, <png file name>
\scalebox{#1}{
	\begin{tikzpicture}
            \begin{axis}[
            	   width = 5.5cm, 
            	   %height = 6.7cm,
                enlargelimits = false,
                axis on top = true,
                axis equal image,
                %unit vector ratio= 0.3 1, % change aspect ratio, one of them should be 1
                point meta min = -1,   
                point meta max = 1,
                xlabel = {$x$ [\SI{}{\milli \metre}]},
                ylabel = {$z$ [\SI{}{\milli \metre}]},
                label style = {font = #2},
                tick label style = {font = #3},
                xlabel style = {yshift = 0.3cm},
                ylabel style = {yshift = -0.3cm},
                y dir = reverse,
                xtick = {15, 20, 25},
                ytick = {15, 20, 25},
                ]
                \addplot graphics [
                    xmin = 15,
                    xmax = 25,
                    ymin = 15,
                    ymax = 25
                ]{#4};
            \end{axis}
	\end{tikzpicture}
	}
}

% cimg for z_def = 30mm
\newcommand{\imgzdefmiddle}[4]{% <scale size>, <label font size>, <tick font size>, <png file name>,
\scalebox{#1}{
	\begin{tikzpicture}
            \begin{axis}[
            	   width = 5.5cm, 
            	   %height = 6.7cm,
                enlargelimits = false,
                axis on top = true,
                axis equal image,
                %unit vector ratio= 0.3 1, % change aspect ratio, one of them should be 1
                point meta min = -1,   
                point meta max = 1,
                xlabel = {$x$ [\SI{}{\milli \metre}]},
                ylabel = {$z$ [\SI{}{\milli \metre}]},
                label style = {font = #2},
                tick label style = {font = #3},
                xlabel style = {yshift = 0.3cm},
                ylabel style = {yshift = -0.3cm},
                y dir = reverse,
                xtick = {15, 20, 25},
                ytick = {25, 30, 35},%21, 22.5, 24
                ]
                \addplot graphics [
                    xmin = 15,
                    xmax = 25,
                    ymin = 25,
                    ymax = 35
                ]{#4};
            \end{axis}
	\end{tikzpicture}
	}
}

% cimg for z_def = 50mm
\newcommand{\imgzdefdeep}[4]{% <scale size>, <label font size>, <tick font size>, <png file name>,
\scalebox{#1}{
	\begin{tikzpicture}
            \begin{axis}[
            	   width = 5.5cm, 
            	   %height = 6.7cm,
                enlargelimits = false,
                axis on top = true,
                axis equal image,
                %unit vector ratio= 0.3 1, % change aspect ratio, one of them should be 1
                point meta min = -1,   
                point meta max = 1,
                xlabel = {$x$ [\SI{}{\milli \metre}]},
                ylabel = {$z$ [\SI{}{\milli \metre}]},
                label style = {font = #2},
                tick label style = {font = #3},
                xlabel style = {yshift = 0.3cm},
                ylabel style = {yshift = -0.3cm},
                y dir = reverse,
                xtick = {15, 20, 25},
                ytick = {45, 50, 55},%21, 22.5, 24
                ]
                \addplot graphics [
                    xmin = 15,
                    xmax = 25,
                    ymin = 45,
                    ymax = 55
                ]{#4};
            \end{axis}
	\end{tikzpicture}
	}
}

% set color

\usepackage{color}
 
\definecolor{fri_gray}{rgb}{0.8, 0.8, 0.8}
\definecolor{fri_green_light}{rgb}{0.76, 0.95, 0.81}
\definecolor{fri_green}{rgb}{0.51, 0.87, 0.78}
\definecolor{tui_orange}{rgb}{0.94, 0.49, 0}
\definecolor{tui_blue}{rgb}{0, 0.22, 0.39}

\definecolor{box_white}{cmyk}{0.0361,0.0251,0.0166,0}
\definecolor{text_black}{cmyk}{0.7979,0.7417,0.6916,0.6554}
\definecolor{tui_orange_dark}{cmyk}{0,0.6,1,0}
\definecolor{tui_orange_light}{cmyk}{0.0000,0.0876, 0.1474, 0.0157}
\definecolor{tui_green_dark}{cmyk}{1,0,0.5,0.2}
\definecolor{tui_green_light}{cmyk}{0.0576,0.0041, 0.0000, 0.0471}
\definecolor{tui_blue_dark}{cmyk}{1.0000,0.5000,0.0000,0.6000}
\definecolor{tui_blue_light}{cmyk}{0.0920,0.0440,0.0000,0.0196}
\definecolor{tui_red_dark}{cmyk}{0.0000,1.0000,1.0000,0.2000}
\definecolor{tui_red_light}{cmyk}{0.0000,0.1107,0.1107,0.0078}

%%%% TikZ setup %%%%

% for SmartInspect
\tikzstyle{scanpath} = [tui_blue, dotted, ->, shorten >=1mm, shorten <=1mm,]
\tikzstyle{scanpoint} = [circle, draw, tui_blue, inner sep = \rCircle, fill = tui_blue]
\tikzstyle{campoint} = [circle, draw, black, inner sep = \rCircleCamera, fill = gray]

% for synthetic aperture
\tikzstyle{griddot} = [circle, draw, black, inner sep = 0.01cm, fill = black]

% for drawing block diagram
\tikzstyle{line} = [draw, thick]
\tikzstyle{line1} = [draw, thick, ->]
\tikzstyle{line3} = [draw, gray, thick, dashed, ->]
\tikzstyle{block} = [rectangle, draw, 
    text width=2cm, text centered, minimum height=0.8cm]
\tikzstyle{block2} = [rectangle, draw, 
    text width=1.8cm, text centered, minimum height=2.8cm]
\tikzstyle{largebox} = [rectangle, draw, dashed, very thick, 
    text width=6.8cm,  minimum height=2cm]
\tikzstyle{block4} = [rectangle, draw, 
    text width=1.5cm, text centered, minimum height=1cm]
\tikzstyle{adder} = [circle, draw, minimum size=2em]
\tikzstyle{junction} = [circle, draw, fill=black, minimum size=1pt]
\tikzstyle{middlepoint} = [circle, draw, fill=gray!30, minimum size=1em]
\tikzstyle{end} = [circle, draw, fill=black!90, minimum size=1em \and circle, draw, minimum size=1.5em]


% for markers in RMSE
\pgfplotsset{
  every axis plot post/.append style={
    every mark/.append style={line width=3pt}
  }
}
% for posscan simulation flow
\tikzstyle{graydashed} = [draw, gray, thick, dashed]

% math operators

\DeclareMathOperator{\sinc}{sinc}
\DeclareMathOperator{\round}{round}
\DeclareMathOperator{\svd}{svd}
\DeclareMathOperator{\argmin}{argmin}
\DeclareMathOperator{\sgn}{sgn}
\DeclareMathOperator{\zeros}{zeros}

\DeclareMathOperator{\dx}{dx}
\DeclareMathOperator{\dy}{dy}
\DeclareMathOperator{\dz}{dz}
\DeclareMathOperator{\dt}{dt}
\DeclareMathOperator{\dist}{d}
\DeclareMathOperator{\pos}{pos}

\DeclareMathOperator{\Expect}{{{\mathbb E}}}

\DeclareMathOperator{\T}{T}

% Metrics
\DeclareMathOperator{\SE}{SE}
\DeclareMathOperator{\MSE}{MSE}
\DeclareMathOperator{\API}{API}
\DeclareMathOperator{\MAPI}{MAPI}
\DeclareMathOperator{\GCNR}{gCNR}
\DeclareMathOperator{\MGCNR}{MgCNR}
\DeclareMathOperator{\OVL}{OVL}

% else
\DeclareMathOperator{\optimized}{opt}
\DeclareMathOperator{\estimated}{est}
\DeclareMathOperator{\thres}{th}
\DeclareMathOperator{\Frob}{F}
\DeclareMathOperator{\Hermit}{H}
\DeclareMathOperator{\TLS}{TLS}
\DeclareMathOperator{\ROI}{ROI}
\DeclareMathOperator{\CF}{CF}





% PDF-options
\title{Iterative SAFT Reconstruction for Manually Acquired Ultrasonic Measurement Data in Nondestructive Testing} 
\subtitle{CSP Research Project 19SS}
\institute{\foreignlanguage{german}{Technische Universität Ilmenau}}
\author{Sayako Kodera}
\date{Aug.07.2019}

% beamer goto-button setup
\setbeamercolor{button}{fg=white,bg=blue}
\renewcommand{\beamergotobutton}[1]{%
    \begingroup% keep color changes local
    \setbeamercolor{button}{fg=white, bg=tui_orange}%
    \beamerbutton{\insertgotosymbol#1}% original definition
    \endgroup
    }

%%%%%%%%%%%%%%%%%%%%%%%%%%%%%%%%%%%%%%%%%%%%%%
%================================= tips to use beamer ======%

% \uncover<page-> = the text will be covered (so displayed in light colors) for the particular pages of the slide
% \ony<page-> = the text/images will be hidden for the particular pages of the slide
% \footnotetext[number] = text for footnote wirh the givien "number"
% \footnotemark = displays the mark (i.e. nuber) of the footnote

%%%%%%%%%%%%%%%%%%%%%%%%%%%%%%%%%%%%%%%%%%%%%%
%%%%%%%%%%%%%%%%%%%%%%%%%%%%%%%%%%%%%%%%%%%%%%

\begin{document}
% display only logos but no contact info in footline
\setbeamertemplate{footline}[light]
% print title slide
\begin{frame}[noframenumbering] % noframenumbering prevents the framecounter from increasing for this single slide
 	\titlepage
\end{frame}

% small footline for more content
\setbeamertemplate{footline}[shrunkplain]
%\begin{frame}[noframenumbering] % noframenumbering prevents the %framecounter from increasing for this single slide
%	\frametitle{Table of contents}
%	\vfill
% 	\tableofcontents[part=1]
%\end{frame}

%%%%%%%%%%%%%%%%%%%%%%%%%%%%%%%%%%%%%%%%%%%%%%
\part{content}

%%%%%%%%%%%%%%%%%%%%%%%%%%%%%%%%%%%%%%%%%%%%%%
%%%%%%%%%%%%%%%%%%%%%%%%%%%%%%%%%%%% Sec. 1.1 %%%%%
\section{Background}

\subsection{Measurement Assistance System} 
% small footline with pagenumbers
\setbeamertemplate{footline}[shrunklight]
% full headline for sections in headline
\setbeamertemplate{frametitle}[full]
\begin{frame}[t]
\frametitle{Measurement Assistance System}
%\fontsize{10pt}
%======================================== content =====%
	\begin{columns}[t]	
	% Text part
	\begin{column}{0.48\textwidth}	
	\vspace{\topsep}
	\hspace*{-0.3cm} \textbf{Features}:
	\vspace*{-0.5cm}
		\begin{itemize}
		\item Position recognition
		\item Data recording
		\item Data visualization
		\item Post-processing
		\end{itemize}
	%	
	\vspace*{0.3cm}
	\textbf{Problem}: Systematic errors\\
	\hspace*{0.3cm} e.g. tracking error 
	\end{column}	
	%
	% Image part
	\begin{column}{0.52\textwidth}
		\begin{overprint} 
		\centering
		\vspace*{-0.5cm}
		\only<1>{
		\inputTikZ{0.6}{figures/SmartInspect_2D.tex}
		}
		\only<2->{
		\inputTikZ{0.6}{figures/tracking_error_2D.tex}
		}
		\end{overprint}
	\end{column}
	\end{columns}	
	
%===================================================%
\end{frame}


\subsection{Impact of Position-Inaccuracy} 
% small footline with pagenumbers
\setbeamertemplate{footline}[shrunklight]
% full headline for sections in headline
\setbeamertemplate{frametitle}[full]
\begin{frame}[t]
\frametitle{Impact of Position-Inaccuracy} 
%\fontsize{10pt}
%======================================== content =====%
	%%=======  Left =======%%%
%	\begin{textblock}{50}(20, 25) %-> using textblock avoid shifting of figures b/w slides
%		\only<1-2>{
%		%\vspace*{10cm}
%		%\centering	
%		\includegraphics[scale=1.2]{images/das_3.pdf}
%		}
%	\end{textblock}
	%
	\begin{textblock}{50}(10, 22)
	\centering
	\begin{overprint}
		% B-Scan ref
		\only<1>{
		\textbf{Correct positions}\\
		\imgbothlabels{1.45}{\scriptsize}{\scriptsize}{figures/pytikz/2D/texpngs/iterative_GD/BScan_190815.png}\\
		% <scale size>, <label font size>, <tick font size>, <png file name>
		}
		% Reco ref
		\only<2>{
		\textbf{Correct positions}\\
		\imgbothlabels{1.45}{\scriptsize}{\scriptsize}{figures/pytikz/2D/texpngs/iterative_GD/Reco_auto_190815.png}\\
		% <scale size>, <label font size>, <tick font size>, <png file name>
		}
	\end{overprint}
	\end{textblock}
	%
	%
	%%======= Right =======%%%
	\begin{textblock}{50}(85, 22)
	\centering
	\begin{overprint}
		%\centering	 -> do not use this! cause horizontal shifting b/w slides!!
%		% B-Scan ref
%		\only<1>{
%		\textbf{Correct positions}\\
%		\imgbothlabels{1.45}{\scriptsize}{\scriptsize}{figures/pytikz/2D/texpngs/iterative_GD/BScan_190815.png}\\
%		% <scale size>, <label font size>, <tick font size>, <png file name>
%		}
%		% B-scan track
%		\only<2>{
%		\textbf{Incorrect positions}\\
%		\imgbothlabels{1.45}{\scriptsize}{\scriptsize}{figures/pytikz/2D/texpngs/iterative_GD/BScan_track_190815_lambda.png}\\
%		% <scale size>, <label font size>, <tick font size>, <png file name>
%		}
		% B-scan track
		\only<1>{
		\textbf{Incorrect positions}\\
		\imgbothlabels{1.45}{\scriptsize}{\scriptsize}{figures/pytikz/2D/texpngs/iterative_GD/BScan_track_190815_lambda.png}\\
		% <scale size>, <label font size>, <tick font size>, <png file name>
		}
		% Reco track
		\only<2>{
		\textbf{Incorrect positions}\\
		\imgbothlabels{1.45}{\scriptsize}{\scriptsize}{figures/pytikz/2D/texpngs/iterative_GD/Reco_track_190815_lambda.png}\\
		% <scale size>, <label font size>, <tick font size>, <png file name>
		}
	\end{overprint}
	\end{textblock}
	
%===================================================%
\end{frame}


\subsection{Objective and Contributions}
% small footline with pagenumbers
\setbeamertemplate{footline}[shrunklight]
% full headline for sections in headline
\setbeamertemplate{frametitle}[full]
\begin{frame}[t]
\frametitle{Objective and Contributions}
%======================================== content =====%
	%%========= Left: Text =========%%
	\begin{textblock}{80}(10, 25) %-> using textblock avoid shifting of figures b/w slides
	\textbf{Objective} 	
	\begin{itemize}
		\item Artefacts reduction
	\end{itemize}
	\vspace*{0.2cm}
	
	\textbf{Contributions}
	\begin{itemize}
		\item Spatial approximation of the reconstruction matrix
		\item Iterative method to correct the positional information
	\end{itemize}
	\end{textblock}
	
	%%========= Right: Image =========%%
	\begin{textblock}{50}(95, 25)
	\centering
	\begin{overprint}
	% Reco track
	\only<1>{
	\imgbothlabels{1.2}{\scriptsize}{\scriptsize}{figures/pytikz/2D/texpngs/iterative_GD/Reco_track_190815_lambda.png}\\
	}
	% Reco opt
	\only<2>{
	\imgbothlabels{1.2}{\scriptsize}{\scriptsize}{figures/pytikz/2D/texpngs/iterative_GD/Reco_opt_190815_lambda.png}\\
	}
	\end{overprint}
	\end{textblock}
	
%===================================================%	
\end{frame}

\section{Method}
\subsection{Iterative Update}
% small footline with pagenumbers
\setbeamertemplate{footline}[shrunklight]
% full headline for sections in headline
\setbeamertemplate{frametitle}[full]
\begin{frame}[t]
\frametitle{Iterative Matrix Improvement} 
%======================================== content =====%
	\textbf{Spatial approximation}
		\begin{itemize}
		\item 1st order Taylor approximation
		\end{itemize}
	%
	\textbf{Iterative position correction}
		\begin{itemize}
		\item Solving the least squares problem
		\item Update the positional information $\xhat$ $\rightarrow$ $\xopt$
		\end{itemize}
	%
	\textbf{Evaluation criterion}
		\begin{itemize}
		\item Error correction $\xdeltaopt = x - \xopt$
		\item Approximation quality $\SEdag =$ $0$ ... $1$
		\end{itemize}
%===================================================%		
\end{frame}


\section{Results}

\subsection{Results: Error Correction}
% small footline with pagenumbers
\setbeamertemplate{footline}[shrunklight]
% full headline for sections in headline
\setbeamertemplate{frametitle}[full]
\begin{frame}[t]
\frametitle{Results: Error Correction} 
%======================================== content =====%
	%%% Positions
	\begin{textblock}{70}(10, 20)
		% Positions
		\only<1-2>{Variables: $|s_{x} - x| $, $\xdelta$\\}
		\only<3>{$|s_{x} - x| = $ \SI{1}{\milli \metre} $\hat{=}$ $0.8 \lambda$\\}
		\only<4>{$|s_{x} - x| = $ \SI{2.5}{\milli \metre} $\hat{=}$ $1.98 \lambda$\\}
		\only<5>{$|s_{x} - x| = $ \SI{5}{\milli \metre}\\}
		%
		\centering
		% Measurement
		%\inputTikZ{1}{figures/measurement_animmation.tex}\\
		\measanimate{1}{2-}{3}{4}{5}\\ 
		% <scale size>, <slide page for base pulse>, <slide page for highlighting the pulse 1mm away>, <slide page for highlighting the pulse 2.5mm away>, <slide page for highlighting the pulse 5mm away>
		%
		\cigray{Wavelength $\lambda$ = \SI{1.26}{\milli \metre}}
	\end{textblock}
	%
	%%% Results
	\begin{textblock}{80}(70, 25)
		\centering
		% errmax_5lambda
		\gdpeanimate{1}{\small}{\normalsize}{figures/pytikz/1D/coordinates/errmax_5lambda/gd_pe_1mm_away.tex}{figures/pytikz/1D/coordinates/errmax_5lambda/gd_pe_2_5mm_away.tex}{figures/pytikz/1D/coordinates/errmax_5lambda/gd_pe_5mm_away.tex}
	\end{textblock}	 
%===================================================%		
\end{frame}


\subsection{Results: SE}
% small footline with pagenumbers
\setbeamertemplate{footline}[shrunklight]
% full headline for sections in headline
\setbeamertemplate{frametitle}[full]
\begin{frame}[t]
\frametitle{Results: A-Scan Approximation} 
%======================================== content =====%
	%%% Positions
	\begin{textblock}{70}(10, 20)
	% Positions
	\only<1>{$|s_{x} - x| = $ \SI{1}{\milli \metre} $\hat{=}$ $0.8 \lambda$\\}
	\only<2>{$|s_{x} - x| = $ \SI{2.5}{\milli \metre} $\hat{=}$ $1.98 \lambda$\\}
	\only<3>{$|s_{x} - x| = $ \SI{5}{\milli \metre}\\}
	%
	\centering
	% Measurement
	%\inputTikZ{1}{figures/measurement_animmation.tex}\\
	\measanimate{1}{1-}{1}{2}{3}\\ 
	% <scale size>, <slide page for base pulse>, <slide page for highlighting the pulse 1mm away>, <slide page for highlighting the pulse 2.5mm away>, <slide page for highlighting the pulse 5mm away>
	%
	\cigray{Wavelength $\lambda$ = \SI{1.26}{\milli \metre}}
	\end{textblock}
	%	
	% Results
	\begin{textblock}{80}(70, 25)
	\centering
	%errmax_5lambda
	\gdseanimate{1}{\small}{\normalsize}{figures/pytikz/1D/coordinates/errmax_5lambda/gd_se_1mm_away.tex}{figures/pytikz/1D/coordinates/errmax_5lambda/gd_se_2_5mm_away.tex}{figures/pytikz/1D/coordinates/errmax_5lambda/gd_se_5mm_away.tex}
	\end{textblock}	 	 
%===================================================%		
\end{frame}


\subsection{Results: Reconstruction} 
% small footline with pagenumbers
\setbeamertemplate{footline}[shrunklight]
% full headline for sections in headline
\setbeamertemplate{frametitle}[full]
\begin{frame}[t]
\frametitle{Results: Reconstruction} 
%\fontsize{10pt}
%======================================== content =====%
	%%=======  Left: B-Scan =======%%%
	\begin{textblock}{75}(10, 22)
	\centering
	\begin{overprint}
		% B-Scan ref
		\only<1>{
		\textbf{No tracking error}\\
		\imgbothlabels{1.45}{\scriptsize}{\scriptsize}{figures/pytikz/2D/texpngs/iterative_GD/BScan_190815.png}\\
		% <scale size>, <label font size>, <tick font size>, <png file name>
		}
		% B-Scan track
		\only<2>{
		\textbf{With $\xdelta = $ \SI{1.26}{\milli \metre}}\\
		\vspace*{0.1cm}
		\imgbothlabels{1.45}{\scriptsize}{\scriptsize}{figures/pytikz/2D/texpngs/iterative_GD/BScan_track_190815_lambda.png}\\
		% <scale size>, <label font size>, <tick font size>, <png file name>
		}
		% B-Scan opt
		\only<3->{
		\textbf{After position correction}\\
		\imgbothlabels{1.45}{\scriptsize}{\scriptsize}{figures/pytikz/2D/texpngs/iterative_GD/BScan_opt_190815_lambda.png}\\
		% <scale size>, <label font size>, <tick font size>, <png file name>
		}
	\end{overprint}
	\end{textblock}
	%
	%
	%%======= Right: Reco =======%%%
	\begin{textblock}{75}(85, 22)
	\centering
	\begin{overprint}
		\textbf{Reconstruction}\\
		% Reco ref
		\only<1>{
		\imgbothlabels{1.45}{\scriptsize}{\scriptsize}{figures/pytikz/2D/texpngs/iterative_GD/Reco_auto_190815.png}\\
		% <scale size>, <label font size>, <tick font size>, <png file name>
		}
		% Reco track
		\only<2>{
		\imgbothlabels{1.45}{\scriptsize}{\scriptsize}{figures/pytikz/2D/texpngs/iterative_GD/Reco_track_190815_lambda.png}\\
		% <scale size>, <label font size>, <tick font size>, <png file name>
		}
		% Reco opt
		\only<3>{
		\imgbothlabels{1.45}{\scriptsize}{\scriptsize}{figures/pytikz/2D/texpngs/iterative_GD/Reco_opt_190815_lambda.png}\\
		% <scale size>, <label font size>, <tick font size>, <png file name>
		}
	\end{overprint}
	\end{textblock}
	
%===================================================%
\end{frame}


%%%%%%%%%%%%%%%%%%%%%%%%%%%%%%%%%%%%%%%%%%%%%%
%%%%%%%%%%%%%%%%%%%%%%%%%%%%%%%%%%%%% Sec. 4 %%%%
\section{Summary} 

\subsection{Conclusion}
% small footline with pagenumbers
\setbeamertemplate{footline}[shrunklight]
% full headline for sections in headline
\setbeamertemplate{frametitle}[full]
\begin{frame}
\frametitle{Conclusion}
%\vspace*{1cm}
%\footnotesize
%======================================== content =====%
	\begin{table}
	\LARGE
	\begin{center}
		\inputTable{0.96\textwidth}{tables/table_conclusion.tex}
	\end{center}
	\end{table}
%===================================================%
\end{frame}


%%%%%%%%%%%%%%%%%%%%%%%%%%%%%%%%%%%%%%%%%%%%%%
%%%%%%%%%%%%%%%%%%%%%%%%%%%%%%%%% Bibliography %%%%%
\bibliography{main}

%%%%%%%%%%%%%%%%%%%%%%%%%%%%%%%%%%%%%%%%%%%%%%
%%%%%%%%%%%%%%%%%%%%%%%%%%%%%%%%%%%% Backup %%%%%
\appendix
\backupbegin
\section{Appendix}

\subsection{Appendix Start}
% small footline with pagenumbers
\setbeamertemplate{footline}[shrunklight]
% full headline for sections in headline
\setbeamertemplate{frametitle}[full]
\begin{frame}
	%\frametitle{ } 
	%\vspace*{1cm}
	%\footnotesize
	%======================================== content =====%
	\centering
	\huge Appendix
	%===================================================%
\end{frame}

%%% Appendix %%%
\subsection{Parameters (1)}
% small footline with pagenumbers
\setbeamertemplate{footline}[shrunklight]
% full headline for sections in headline
\setbeamertemplate{frametitle}[full]
\begin{frame}
	\frametitle{Parameters w.r.t. Test Object}
	%\vspace*{1cm}
	%\footnotesize
	%======================================== content =====%
	\centering
	\begin{table}
	\begin{center}
		\inputTable{0.85\textwidth}{tables/table_constant_params_EN.tex}
	\end{center}
	\end{table}
	%===================================================%
\end{frame}


\subsection{Parameters (2)}
% small footline with pagenumbers
\setbeamertemplate{footline}[shrunklight]
% full headline for sections in headline
\setbeamertemplate{frametitle}[full]
\begin{frame}
	\frametitle{Parameters w.r.t. Pulse}
	%\vspace*{1cm}
	%\footnotesize
	%======================================== content =====%
	\centering
	\begin{columns}
	% table
	\begin{column}{0.5\textwidth}
	\begin{table}
	\begin{center}
		\inputTable{1\textwidth}{tables/table_constant_params_pulse_EN.tex}
	\end{center}
	\end{table}
	\end{column}
	% FIG : Gabot
	\begin{column}{0.5\textwidth}
	\inputTikZ{0.6}{figures/Gabor_pulse.tex}
	\end{column}
	\end{columns}	
	%===================================================%
\end{frame}


%%%%%%%%%%%%%%%%%%%%%%%%%%%%%%%%%%%%%%%%%%%%%%
%%%%%%%%%%%%%%%%%%%%%%%%%%%%%%%%%%% Sec. 1.2 %%%%%


\subsection{SAFT}
% small footline with pagenumbers
\setbeamertemplate{footline}[shrunklight]
% full headline for sections in headline
\setbeamertemplate{frametitle}[full]
\begin{frame}[t]
\frametitle{Post-processing Method} 
%======================================== content =====%
	Synthetic Aperture Focusing Technique (SAFT)
	\begin{columns}[t]
	% Text part
	\begin{column}{0.48\textwidth}
	\begin{center}
		\begin{itemize}
		\item Superposition according to propagation time delay
		\item Spatial sampling of the specimen
		\end{itemize}
	\end{center}
	\end{column}	
	%	
	% Image part
	\begin{column}{0.52\textwidth}
	\vspace*{-0.3cm}
	\begin{center}
		\only<1>{\inputTikZ{1.1}{figures/SAFT_superposition.tex}}
		\only<2>{\inputTikZ{0.9}{figures/Synthetic_aperture.tex}}
	\end{center}
	\end{column}
	%
	\end{columns}	
%===================================================%		
\end{frame}

\subsection{FWM}
% small footline with pagenumbers
\setbeamertemplate{footline}[shrunklight]
% full headline for sections in headline
\setbeamertemplate{frametitle}[full]
\begin{frame}[t]
\frametitle{Mathematical Model} 
%======================================== content =====%
	% Defect map
	\only<1>{
	\textit{defect map} for ROI = \cihighlight{$\M \times \N_{x}$}\\
	\centering
	\inputTikZ{1.1}{figures/ROI_sketch.tex}
	}
	%
	% Linear transform: single A-Scan
	\only<2-3>{
	\textbf{Transform}: \textit{defect map} $\rightarrow$ single A-Scan at $x_{k}$ $\in$ $\RR$\\
	% Equation
	\begin{equation*}
	\ascank = \SAFTk \defect + \noisevec \in \RR^{\M}
	\end{equation*}\\ % "\\" avoids to shift the next text horizontaly
	%
	\vspace*{0.2cm}
	% Description
	\only<2>{
		\begin{itemize}
			\item $\ascank$: measured A-Scan at $x_{k}$ (\cihighlight{$\M$})
			\item $\SAFTk$: SAFT matrix at $x_k$ (\cihighlight{$\M \times \LL = \M \times \M \N_{x}$})\\
					\hspace*{0.7cm} Containing pulse information for $\scatterer_{l} = (x_i, z_j)$, $l = 1$...$\LL$
			\item $\defect$: vectorized \textit{defect map} (\cihighlight{$\LL = \M \N_{x}$})
			\item $\noisevec$: noise (\cihighlight{$\M$})
		\end{itemize}
	}
	\only<3>{
		% SAFT_k
		\textbf{SAFT matrix}  at $x_k$\\
		\begin{equation*}
			\SAFTk = 
			\begin{bmatrix} \SAFTcol{k}{1} & \SAFTcol{k}{2} & \cdots & \SAFTcol{k}{\LL} \end{bmatrix}
			\in \RR^{\M \times \LL}
		\end{equation*}
	}}
	%
	% Linear transform: K A-Scans
	\only<4->{
	\textbf{Transform}: \textit{defect map} $\rightarrow$ $\K$ A-Scans at $\xvec$ $\in$ $\RR^{\K}$\\
	% Equation
	\begin{equation*}
	\vecoperator{\Ascan} = \SAFT \defect + \vecoperator{\bm{N}} \in \RR^{\M \K}
	\end{equation*}\\ % "\\" avoids to shift the next text horizontaly
	%
	\vspace*{0.2cm}
	% Description
	\only<4>{
	\begin{itemize}
	\item $\Ascan$: measured A-Scans at $\xvec$ (\cihighlight{$\M \times \K$})
	\item $\SAFT$: SAFT matrix at $\xvec$ (\cihighlight{$\M \K \times \LL = \M \times \M \N_{x}$})\\
			\hspace*{0.7cm} Containing pulse information for $\scatterer_{l} = (x_i, z_j)$, $l = 1$...$\LL$
	\item $\defect$: vectorized \textit{defect map} (\cihighlight{$\LL = \M \N_{x}$})
	\item $\bm{N}$: noise (\cihighlight{$\M \times \K$})
	\end{itemize}
	}
	\only<5>{
		% A-Scans
		\textbf{A-Scans}  collected at $\xvec$\\
		\begin{equation*}
			\Ascan = 
		\begin{bmatrix} \ascanvec_{1} & \ascanvec_{2} & \cdots & \ascanvec_{\K} \end{bmatrix}
			\in \RR^{\M \times \K}
		\end{equation*}
	}
	\only<6>{
		% SAFT
		\textbf{SAFT matrix}  at $\xvec$\\
		\begin{equation*}
			\SAFT = 
			\begin{bmatrix} \SAFT_1 \\ \SAFT_2 \\ \vdots \\ \SAFT_{\K} \end{bmatrix}
			\in \RR^{\M \K \times \LL}
		\end{equation*}
	}}
	%
%	% Matrix description
%	\only<6>{
%	% SAFT_k
%	\begin{equation*}
%		\SAFTk = 
%		\begin{bmatrix} \SAFTcol{k}{1} & \SAFTcol{k}{2} & \cdots & \SAFTcol{k}{\LL} \end{bmatrix}
%	\end{equation*}
%	%
%	% A-Scans
%	\begin{equation*}
%		\Ascan = 
%		\begin{bmatrix} \ascanvec_{1} & \ascanvec_{2} & \cdots & \ascanvec_{\K} \end{bmatrix}
%	\end{equation*}
%	%
%	% SAFT matrix
%	\begin{equation*} 
%		\SAFT = \begin{bmatrix} \SAFT_1 \\ \SAFT_2 \\ \vdots \\ \SAFT_{\K} \end{bmatrix}.
%	\end{equation*}
%	}
%===================================================%		
\end{frame}


\subsection{Spatial Approximation}
% small footline with pagenumbers
\setbeamertemplate{footline}[shrunklight]
% full headline for sections in headline
\setbeamertemplate{frametitle}[full]
\begin{frame}[t]
\frametitle{Spatial Approximation} 
%======================================== content =====%
	% Single A-Scan
	\only<1>{
	\textbf{\cihighlight{Tracked position}}: $\xhat_{k}$ $=$ $x_k + \Delta x_k$ $\in$ $\RR$ \\
	\vspace*{0.3cm}
	%
	\hspace*{0.5cm} \textbf{Data model}\\
	\vspace*{-0.5cm}
	\begin{equation*}
	\ascank = \SAFTk \defect + \noisevec \in \RR^{\M}
	\end{equation*}
	%
	\hspace*{0.5cm} \textbf{Approximation}\\
	\vspace*{-0.5cm}
	\begin{equation*}
	\SAFTk \approx \SAFThatk - \Jacobianhat_{k} \Delta x_k \in \RR^{\M \times \LL}
	\end{equation*}
	%
	% Description
	\vspace*{-0.5cm}
	\begin{itemize}
	\item $\SAFThatk$: SAFT matrix at $\xhat_{k}$ (\cihighlight{$\M \times \LL$})
	\item $\Jacobianhat_{k}$: Derivative of $\SAFThatk$ w.r.t. $x$ (\cihighlight{$\M \times \LL$})
	\end{itemize}
	%
	}
	%
	% multiple (K) A-Scans
	\only<2->{ 
	\textbf{\cihighlight{Tracked positions}}: $\xvechat$ $=$ $\xvec + \Delta \xvec$ $\in$ $\RR^{\K}$\\
	\vspace*{0.3cm}
	%
	\hspace*{0.5cm} \textbf{Data model}\\
	\vspace*{-0.5cm}
	\begin{equation*}
	\vecoperator{\Ascan} = \SAFT \defect + \vecoperator{\bm{N}} \in \RR^{\M \K}
	\end{equation*}
	%
	\hspace*{0.5cm} \textbf{Approximation}\\
	\vspace*{-0.5cm}
	\begin{equation*}
	\SAFT \approx \SAFThat + \bm{E} \Jacobianhat \in \RR^{\M \K \times \LL}
	\end{equation*}
	%
	% Description
	\vspace*{-0.5cm}
	\begin{itemize}
	\item $\SAFThat$: SAFT matrix at $\xvechat$ (\cihighlight{$\M \K \times \LL$})
	\item $\Jacobianhat$: Derivative of $\SAFThat$ w.r.t. $\xvec$ (\cihighlight{$\M \K \times \LL$})
	\item $\bm{E}$: Error matrix $= \diagoperator{\Delta \xvec} \otimes \Identity_{\M}$ (\cihighlight{$\M \K \times \M \K$})
	\end{itemize}
	}
%===================================================%		
\end{frame}


\subsection{Proposed Preprocessing}
% small footline with pagenumbers
\setbeamertemplate{footline}[shrunklight]
% full headline for sections in headline
\setbeamertemplate{frametitle}[full]
\begin{frame}[t]
\frametitle{Proposed Preprocessing} 
%======================================== content =====%
	% Cost function
	\begin{equation*}
	\min_{\xdeltavec} \min_{\defect} \norm{ \vecoperator{\Ascan}  - \SAFTapproxxvec \defect }_{2}^{2}
	\end{equation*}
	%
	% Block diagram
	\begin{figure}
	\inputTikZ{1}{figures/block_diagram_simulation_flow_base.tex}
	\end{figure}
	%
%===================================================%		
\end{frame}


\subsection{Curve Fitting}
% small footline with pagenumbers
\setbeamertemplate{footline}[shrunklight]
% full headline for sections in headline
\setbeamertemplate{frametitle}[full]
\begin{frame}[t]
\frametitle{Curve Fitting} 
%======================================== content =====%
	\only<1>{%
		% Cost function
		%\vspace*{-0.18cm} % to compensate the vertical shift
		\begin{equation*}
		\min_{\xdeltavec} \cihighlight{\min_{\defect}} \norm{ \vecoperator{\Ascan}  - \SAFTapproxxvec \cihighlight{\defect}}_{2}^{2}
		\end{equation*}
		%
		% Block diagram
		\begin{figure}
		\inputTikZ{1}{figures/block_diagram_simulation_flow_hyperbolafit.tex}
		\end{figure}
		%
	}%
	\only<2->{%
	\begin{textblock}{80}(10, 25)
		\only<2-4>{%
			\textbf{B-Scan}
			\begin{itemize}
				\item Defect map $\defect$ $\hat{=}$ \cihighlight{$(x_{\dist}, z_{\dist})$} 
				\item Scan positions $\xvec$ $\in$ $\RR^{\K}$
				\item Peak positions $\zvec$ $\in$ $\RR^{\K}$ \\
						$\hat{=}$ Time-of-Flight
				\item Curvature
			\end{itemize}
		$\Rightarrow$ \textbf{$\approx$ Parabola}
		}%
		%
		\only<5->{%
			% Eq: parabola approx. polynomial form
			\textbf{Parabola approximation}: $k$-th scan\\ \vspace*{-0.5cm}
			\begin{equation*}
			z_{k} \approx  u_{0} + u_1 \cdot x_{k} + u_{2} \cdot x_{k}^{2}
			\end{equation*}
			%
			\only<5>{%
				$\Rightarrow$ \hspace*{1.4cm} $x_{\dist} = - \frac{u_{1}}{2 u_{2}}$\\
				\vspace*{0.2cm}
				\hspace*{2.2cm}$z_{\dist} = u_{0} - \frac{u_{1}^{2}}{4 u_{2}}$\\
			}%
			\only<6->{
				% K scan positions
				\textbf{$\K$ scan positions}:
				\begin{equation*}
				\begin{bmatrix} z_{1}\\ z_{2}\\  \vdots \\ z_{\K} \end{bmatrix} \approx 
				\begin{bmatrix} 1 & x_{1} & x_{1}^{2}\\ 1 & x_{2} & x_{2}^{2}\\ \vdots & \vdots & \vdots \\ 1 & x_{\K} & x_{\K}^{2} \end{bmatrix} \cdot
				\begin{bmatrix} u_{0}\\ u_{1}\\ u_{2} \end{bmatrix}
				\end{equation*}
				%
				\hspace*{1.4cm} $\bm{z}$ \hspace*{1.7cm} $\bm{X}$ \hspace*{1.5cm} $\bm{u}$
			}
		}%
		%
	\end{textblock}
	%
	%
	\begin{textblock}{80}(90, 27.5)
		\centering
		% Measurement
		%\inputTikZ{1}{figures/measurement_animmation.tex}\\
		\curvefit{1}{4-}{5-}{3-}{5-}\\ 
		% <scale size>, <slide page for base pulse>, <slide page for highlighting the pulse 2.5mm away>, <slide page for adding the defect position>, <slide page for adding the scan position>, 
		%
	\end{textblock}	 
	}
%===================================================%		
\end{frame}


\subsection{TLS}
% small footline with pagenumbers
\setbeamertemplate{footline}[shrunklight]
% full headline for sections in headline
\setbeamertemplate{frametitle}[full]
\begin{frame}[t]
\frametitle{Polynomial Regression via TLS} 
%======================================== content =====%
	\only<1>{%
	% Data approximation: K scan positions
	\begin{textblock}{80}(10, 25)
		\textbf{Approximation}
		\begin{equation*}
		\begin{bmatrix} z_{1}^{2}\\ z_{2}^{2}\\  \vdots \\ z_{\K}^{2} \end{bmatrix} \approx 
		\begin{bmatrix} 1 & x_{1} & x_{1}^{2}\\ 1 & x_{2} & x_{2}^{2}\\ \vdots & \vdots & \vdots \\ 1 & x_{\K} & x_{\K}^{2} \end{bmatrix} \cdot
		\begin{bmatrix} u_{0}\\ u_{1}\\ u_{2} \end{bmatrix}
		\end{equation*}
		%
		\hspace*{1.4cm} $\bm{z}$ \hspace*{1.7cm} $\bm{X}$ \hspace*{1.5cm} $\bm{u}$\\
	\end{textblock}	
	%
	% Errors
	\begin{textblock}{60}(85, 30)
		Errors in $\bm{z}$ and $\bm{X}$\\
		\begin{itemize}
			\item $\Delta \bm{z}$: Measurement noise, quantization error etc
			\item $\Delta \bm{X}$:Tracking error
		\end{itemize}
		$\Rightarrow$ \textbf{independent errors}
	\end{textblock}	
	}%
	%
	\only<2->{%
	\begin{textblock}{150}(10, 25)
		\only<2>{\textbf{Approximation}\\}%
		\only<3->{\textbf{Incorporating errors}\\}%
		\vspace*{0.3cm}
		% Equation
		\only<2>{%
			\hspace*{4.3cm} $\bm{z}$ \hspace*{0.8cm} $\approx$ \hspace*{0.8cm} $\bm{X}$ \hspace*{0.8cm} $\cdot$ \hspace*{0.8cm} $\bm{u}$\\%
		}%
		\only<3->{%
			\hspace*{3.3cm} $\bm{z} + \cihighlight{\Delta \bm{z}}$ \hspace*{0.4cm} $=$ \hspace*{0.4cm} $\left( \bm{X} + \cihighlight{\Delta \bm{X}} \right)$ \hspace*{0.2cm} $\cdot$ \hspace*{0.2cm} $\bm{u}$\\%
		}%
		\vspace*{0.3cm}
		% Dimension
		\cigray{\hspace*{4.3cm} \cigray{$\K$} \hspace*{2cm} $\K \times 3$ \hspace*{1.8cm} $3$}\\
	\end{textblock}	
	%
	\begin{textblock}{150}(10, 50)
		\only<4->{
			$\Rightarrow$ \textbf{Total least squares}\\
			\vspace*{-0.7cm}
				\begin{eqnarray*}
					\min \norm{\left[ \Delta \bm{X} \cigray{\vert} \Delta \bm{z} \right] }_{\F}^{2} & 
					\mbox{s.t.} & 
					\bm{z} + \Delta \bm{z} = \left[ \bm{X} + \Delta \bm{X} \right] \cdot \bm{u} \\
				\end{eqnarray*}
		}
	\end{textblock}	
	%
	\begin{textblock}{150}(10, 70)
		\only<4->{
			\hspace*{1cm} $\Rightarrow$ $\bm{u}$ $\Rightarrow$ $\defecthat_{\estimated}$
		}
	\end{textblock}	
	}%
	%
%===================================================%		
\end{frame}

\subsection{Iterative Error Correction}
% small footline with pagenumbers
\setbeamertemplate{footline}[shrunklight]
% full headline for sections in headline
\setbeamertemplate{frametitle}[full]
\begin{frame}[t]
\frametitle{Iterative Error Correction} 
%======================================== content =====%
	\only<1>{
	% Cost function
	\begin{equation*}
	\cihighlight{\min_{\xdeltavec}} \min_{\defect} \norm{ \vecoperator{\Ascan}  - \left( \SAFThat + \cihighlight{\bm{E} \Jacobianhat} \right) \defect }_{2}^{2}
	\end{equation*}
	%
	% Block diagram
	\begin{figure}
	\inputTikZ{1}{figures/block_diagram_simulation_flow_Newton.tex}
	\end{figure}
	}
	\only<2->{%
	\begin{textblock}{150}(10, 25)
		\textbf{Taylor approximation}
		\begin{equation*}
		\vecoperator{\Ascan} \approx \SAFTapproxxvec \defect 
		\end{equation*}
		%
		\textbf{Cost function (2nd step)}
		\only<2->{%
		\begin{equation*}
		\min_{\xdeltavec} \norm{ \vecoperator{\Ascan} - \SAFTapproxxvec \cihighlight{\defecthat_{\estimated}} }_{2}^{2}
		\end{equation*}
		}%
	\end{textblock}	
	%
	\begin{textblock}{150}(10, 75)
		\only<3->{%
		$\Rightarrow$ Nonlinear programming (e.g. Newton method)
		}%
	\end{textblock}	
	}%
	%
%===================================================%		
\end{frame}


\subsection{TLS Performance}
% small footline with pagenumbers
\setbeamertemplate{footline}[shrunklight]
% full headline for sections in headline
\setbeamertemplate{frametitle}[full]
\begin{frame}[t]
\frametitle{Signal Source Estimation} 
%======================================== content =====%
	% SE result
	\begin{textblock}{80}(10, 25)
		\textbf{Estimation error}\\%
		%
		\meTLS{0.8}{\large}{\large}{1-}{figures/pytikz/1D/TLS_error_762dz_200207_lambda.tex}{2-}{figures/pytikz/1D/TLS_error_1270dz_200207_lambda.tex} 
		% <scale size>, <label font size>, <tick font size>,  <slide page to pop up the fname for ME 762dz>, <fname for ME 762dz>, <slide page to pop up the ME 1270dz>, <fname for ME 1270dz>
		
	\end{textblock}
	%
	% TikZ image
	\begin{textblock}{80}(90, 25)
		\centering
		\TLSanimate{0.9}{1}{2}
		% <scale size>,  <slide page for z_def = 762dz>, <slide page for z_def = 1270dz>
	\end{textblock}	 
%===================================================%		
\end{frame}

\subsection{Metrics: $\SEdag$} 
% small footline with pagenumbers
\setbeamertemplate{footline}[shrunklight]
% full headline for sections in headline
\setbeamertemplate{frametitle}[full]
\begin{frame}
\frametitle{Metrics: $\SEdag$}
%\vspace*{1cm}
%\footnotesize
%======================================== content =====%
	\begin{equation*}
	\SEdag = \frac{ \| \gamma \ascanvechat - \ascanvec \|_{2}}{\| \ascanvec \|_{2}}
	\end{equation*}\\
	%
	\vspace*{0.5cm}
	$\gamma$: normalization factor \\
	\begin{equation*}
	\gamma = \frac{\ascanvec^{\T} \cdot \ascanvechat}{\ascanvechat^{\T} \cdot \ascanvechat}.
	\end{equation*}
%===================================================%
\end{frame}


\subsection{Metrics: $\API$}
% small footline with pagenumbers
\setbeamertemplate{footline}[shrunklight]
% full headline for sections in headline
\setbeamertemplate{frametitle}[full]
\begin{frame}[t]
\frametitle{Metrics: $\API$}
%======================================== content =====%
	\begin{equation*}
		\API = \frac{A_{\epsilon}}{\lambda^{2}} 
	\end{equation*}
%===================================================%	
\end{frame}


\subsection{Metrics: $\GCNR$}
% small footline with pagenumbers
\setbeamertemplate{footline}[shrunklight]
% full headline for sections in headline
\setbeamertemplate{frametitle}[full]
\begin{frame}[t]
\frametitle{Metrics: $\GCNR$}
%======================================== content =====%
	\begin{equation*}
		\OVL = \int \min \{ p_i (x), p_o (x) \} \dx
	\end{equation*}
	where
	\begin{itemize}
		\item $p_i (x)$: p.d.f for inside the target area
		\item $p_o (x)$: p.d.f for outside the target area
	\end{itemize}
	%
	\vspace*{0.5cm}
	\begin{equation*}
		\GCNR = 1 - \OVL
	\end{equation*}
%===================================================%	
\end{frame}

\subsection{Error tolerance Images}
% small footline with pagenumbers
\setbeamertemplate{footline}[shrunklight]
% full headline for sections in headline
\setbeamertemplate{frametitle}[full]
\begin{frame}[t]
\frametitle{Error Tolerance: Images}
%======================================== content =====%
	\vspace*{-0.5cm}
	\only<1>{\textbf{Error = $0.4 \lambda$}}\only<2>{\textbf{Error = $0.8 \lambda$}}\only<3>{\textbf{Error = $1 \lambda$}} ($\lambda = $ \SI{1.26}{\milli \metre})\\
	%%=======  Left: Reco true =======%%%
	\begin{textblock}{50}(5, 27)
	\centering
	\begin{overprint}
		%% 30mm
			\textbf{Reference}\\
			\imgzdefmiddle{1.2}{\scriptsize}{\scriptsize}{figures/pytikz/2D/texpngs/ARP/30mm_true.png}\\
			% <scale size>, <label font size>, <tick font size>, <png file name>
	\end{overprint}
	\end{textblock}
	%
	%
	%%======= Middle: Reco track =======%%%
	\begin{textblock}{50}(55, 27)
	\centering
	\begin{overprint}
		%% 30mm 0.4 lambda
		\only<1>{
			\textbf{No error correction}\\
			\imgzdefmiddle{1.2}{\scriptsize}{\scriptsize}{figures/pytikz/2D/texpngs/ARP/30mm_track_04lambda.png}\\
			% <scale size>, <label font size>, <tick font size>, <png file name>
		}
		%% 30mm 0.8 lambda
		\only<2>{
			\textbf{No error correction}\\
			\imgzdefmiddle{1.2}{\scriptsize}{\scriptsize}{figures/pytikz/2D/texpngs/ARP/30mm_track_08lambda.png}\\
			% <scale size>, <label font size>, <tick font size>, <png file name>
		}
		%% 30mm 1 lambda
		\only<3>{
			\textbf{No error correction}\\
			\imgzdefmiddle{1.2}{\scriptsize}{\scriptsize}{figures/pytikz/2D/texpngs/ARP/50mm_track_1lambda.png}\\
			% <scale size>, <label font size>, <tick font size>, <png file name>
		}
	\end{overprint}
	\end{textblock}
	%
	%
	%%======= Right: Reco opt =======%%%
	\begin{textblock}{50}(105, 27)
	\centering
	\begin{overprint}
		%% 30mm 0.4 lambda
		\only<1>{
			\textbf{With BEC}\\
			\imgzdefmiddle{1.2}{\scriptsize}{\scriptsize}{figures/pytikz/2D/texpngs/ARP/30mm_opt_04lambda.png}\\
			% <scale size>, <label font size>, <tick font size>, <png file name>
		}
		%% 30mm 0.8 lambda
		\only<2>{
			\textbf{With BEC}\\
			\imgzdefmiddle{1.2}{\scriptsize}{\scriptsize}{figures/pytikz/2D/texpngs/ARP/30mm_opt_08lambda.png}\\
			% <scale size>, <label font size>, <tick font size>, <png file name>
		}
		%% 30mm 1 lambda
		\only<3>{
			\textbf{With BEC}\\
			\imgzdefmiddle{1.2}{\scriptsize}{\scriptsize}{figures/pytikz/2D/texpngs/ARP/30mm_opt_1lambda.png}\\
			% <scale size>, <label font size>, <tick font size>, <png file name>
		}
	\end{overprint}
	\end{textblock}	
%===================================================%	
\end{frame}

\subsection{Error tolerance: $\SEdag$}
% small footline with pagenumbers
\setbeamertemplate{footline}[shrunklight]
% full headline for sections in headline
\setbeamertemplate{frametitle}[full]
\begin{frame}[t]
\frametitle{Error tolerance Results: $\SEdag$}
%======================================== content =====%
	\resultSE{1.55}{\scriptsize}{Tracking error / $\lambda$}{0.6}{0, 0.2, ..., 1.01}{figures/pytikz/1D/se_track_tolerance_30mm.tex}{figures/pytikz/1D/se_opt_tolerance_30mm.tex} 
%===================================================%	
\end{frame}

\subsection{Error tolerance: $\API$}
% small footline with pagenumbers
\setbeamertemplate{footline}[shrunklight]
% full headline for sections in headline
\setbeamertemplate{frametitle}[full]
\begin{frame}[t]
\frametitle{Error tolerance Results: $\API$}
%======================================== content =====%
	\resultAPI{1.55}{\scriptsize}{Tracking error / $\lambda$}{26.5}{0, 0.2, ..., 1.01}{18, 20, ..., 26}{figures/pytikz/1D/api_true_tolerance_30mm.tex}{figures/pytikz/1D/api_track_tolerance_30mm.tex}{figures/pytikz/1D/api_opt_tolerance_30mm.tex} 
%===================================================%	
\end{frame}


\subsection{Error tolerance: $\GCNR$}
% small footline with pagenumbers
\setbeamertemplate{footline}[shrunklight]
% full headline for sections in headline
\setbeamertemplate{frametitle}[full]
\begin{frame}[t]
\frametitle{Error tolerance Results: $\GCNR$}
%======================================== content =====%
	\resultGCNR{1.55}{\scriptsize}{Tracking error / $\lambda$}{0.85}{0, 0.2, ..., 1.01}{0.85, 0.9, 0.95}{figures/pytikz/1D/gcnr_true_tolerance_30mm.tex}{figures/pytikz/1D/gcnr_track_tolerance_30mm.tex}{figures/pytikz/1D/gcnr_opt_tolerance_30mm.tex}
%===================================================%	
\end{frame}


\subsection{Impact of the ROI Depth: $\SEdag$}
% small footline with pagenumbers
\setbeamertemplate{footline}[shrunklight]
% full headline for sections in headline
\setbeamertemplate{frametitle}[full]
\begin{frame}[t]
\frametitle{Impact of the ROI Depth: $\SEdag$}
%======================================== content =====%
	\resultSE{1.55}{\scriptsize}{ROI depth [\SI{}{\milli \metre}]}{0.8}{20, 30, ..., 80}{figures/pytikz/1D/se_track_depth.tex}{figures/pytikz/1D/se_opt_depth.tex}
%===================================================%	
\end{frame}


\subsection{Impact of the ROI Depth: $\API$}
% small footline with pagenumbers
\setbeamertemplate{footline}[shrunklight]
% full headline for sections in headline
\setbeamertemplate{frametitle}[full]
\begin{frame}[t]
\frametitle{Impact of the ROI Depth: $\API$}
%======================================== content =====%
	\resultAPI{1.55}{\scriptsize}{ROI depth [\SI{}{\milli \metre}]}{28}{20, 30, ..., 80}{18, 20, ..., 26}{figures/pytikz/1D/api_true_depth.tex}{figures/pytikz/1D/api_track_depth.tex}{figures/pytikz/1D/api_opt_depth.tex} 
%===================================================%	
\end{frame}


\subsection{Impact of the ROI Depth: $\GCNR$}
% small footline with pagenumbers
\setbeamertemplate{footline}[shrunklight]
% full headline for sections in headline
\setbeamertemplate{frametitle}[full]
\begin{frame}[t]
\frametitle{Impact of the ROI Depth: $\GCNR$}
%======================================== content =====%
	\resultGCNR{1.55}{\scriptsize}{ROI depth [\SI{}{\milli \metre}]}{0.65}{20, 30, ..., 80}{0.65, 0.75, ..., 1.01}{figures/pytikz/1D/gcnr_true_depth.tex}{figures/pytikz/1D/gcnr_track_depth.tex}{figures/pytikz/1D/gcnr_opt_depth.tex}	
%===================================================%	
\end{frame} 

\backupend
%%%%%%%%%%%%%%%%%%%%%%%%%%%%%%%%%%%%%%%%%%%%%%
\end{document}
